% =====================================================================
%  4.3  Atividade 3.2.2.3 -- Sistemas coletivos de abastecimento de água
% =====================================================================
\subsection{Atividade 3.2.2.3 -- Implantar e reabilitar sistemas coletivos de abastecimento de água em comunidades rurais}
\label{subsec:abastecimento}

\subsubsection{Objetivo e fundamentação}

Implantar ou reabilitar sistemas coletivos de abastecimento de água, com
captação subterrânea por poço tubular, em 83 comunidades rurais, atendendo
5.810 domicílios. Os procedimentos consideram que as comunidades serão
definidas pelo trabalho técnico-social (\atividade{3.2.2.2}) e que o modelo
de gestão será pactuado com as comunidades antes da instalação dos sistemas
(\atividade{3.2.2.1}) \cite{mop2026}.

\begin{fichaatividade}{Ficha-síntese da Atividade 3.2.2.3}{qua:ficha-3223}
  \campo{Código}{3.2.2.3}
  \campo{Tipo de atividade}{Serviço e obra}
  \campo{Forma de execução}{Contratação pública (licitação), com contratação integrada}
  \campo{Fonte de recursos}{BIRD}
  \campo{Responsável}{IAT / IDR-Paraná}
  \campo{Dependências institucionais}{Prefeituras}
  \campo{Prazo estimado}{Ano 2 ao ano 5}
  \campo{Produto esperado}{83 sistemas implantados}
  \campo{Unidade de medida}{Sistema implantado}
  \campo{Evidência de comprovação}{Termo de Recebimento Definitivo (TRD)}
  \campo{Periodicidade}{Por lote}
  \campo{Validação}{IAT}
\end{fichaatividade}

\subsubsection{Procedimentos}

\paragraph{Fase 1 -- Estudos preliminares e captação.}
\begin{enumerate}[label=\alph*)]
  \item \textbf{estudos preliminares e anteprojeto}: avaliação da viabilidade
        técnica, social e ambiental do sistema; identificação da titularidade
        das áreas de perfuração e de reservatório; verificação da situação
        fundiária das áreas das estruturas;
  \item \textbf{levantamento de campo}: estudos geológicos e
        hidrogeológicos para identificar fontes de água, características do
        solo, aquíferos e condições ambientais;
  \item \textbf{locação do ponto de captação}: definição do local mais
        adequado para a captação e a perfuração do poço;
  \item \textbf{contratação integrada do poço}: a mesma empresa ou consórcio
        será responsável pelo projeto executivo e pela execução da obra;
  \item \textbf{teste de vazão e potabilidade}: verificação da quantidade e
        da qualidade da água disponível;
  \item \textbf{estudos ambientais e sociais para obtenção da outorga}.
\end{enumerate}

\paragraph{Fase 2 -- Projeto e licenciamento.}
\begin{enumerate}[label=\alph*), resume]
  \item \textbf{projeto conceitual ou anteprojeto da solução} de tratamento
        e armazenamento;
  \item \textbf{ordenação de despesas}: autorização administrativa das
        despesas;
  \item \textbf{contratação integrada} do sistema de tratamento,
        reservação e distribuição;
  \item \textbf{projeto básico}: detalhamento técnico suficiente para a
        licitação da obra;
  \item \textbf{projeto executivo}: projeto completo, com todos os detalhes
        necessários à execução da obra;
  \item \textbf{estudos ambientais e sociais para o licenciamento};
  \item \textbf{obtenção das licenças} ambientais e autorizações necessárias.
\end{enumerate}

\paragraph{Fase 3 -- Execução, supervisão e fiscalização.}
\begin{enumerate}[label=\alph*), resume]
  \item \textbf{contratação da supervisão e do gerenciamento da obra}:
        equipe técnica independente para acompanhar a execução;
  \item \textbf{execução da obra}: construção da infraestrutura do sistema;
  \item \textbf{supervisão e gerenciamento da obra}, com \textbf{relatórios
        de supervisão} que registram o andamento, os problemas e as
        recomendações;
  \item \textbf{fiscalização da obra}, com \textbf{relatórios de
        fiscalização} que atestam o atendimento aos projetos e às normas.
\end{enumerate}

\paragraph{Fase 4 -- Entrega e encerramento.}
\begin{enumerate}[label=\alph*), resume]
  \item \textbf{entrega do sistema pelo IAT à entidade gestora}:
        transferência da responsabilidade pela operação e manutenção;
  \item \textbf{encerramento do processo}: conclusão das etapas
        administrativas e técnicas.
\end{enumerate}

\pendencia{O MOP define ``contratação integrada'' como aquela em que o
contratado elabora o projeto executivo e executa a obra. Na Lei
nº~14.133/2021 (art.~6º, XXXII e XXXIII), essa definição corresponde à
contratação \textit{semi-integrada}; na integrada, o contratado elabora também
o projeto básico. Conferir o regime pretendido e a compatibilidade com o
Regulamento de Aquisições do Banco Mundial. A sequência ``contratação
integrada $\rightarrow$ projeto básico'' da Fase~2 sugere o regime
integrado.}

\subsubsection{Etapas, produtos e evidências}

\begin{landscape}
\begin{quadro}[p]
\caption{Etapas de implantação dos sistemas coletivos de abastecimento de água (\atividade{3.2.2.3})}
\label{qua:etapas-3223}
\scriptsize
\renewcommand{\arraystretch}{1.25}
\begin{tabularx}{\linewidth}{|P{2.6cm}|L|P{2.1cm}|Q{1.3cm}|P{2.8cm}|P{2cm}|P{2.4cm}|Q{1.4cm}|P{1.7cm}|}
\hline
\cab{Etapa} & \cab{Descrição técnica} & \cab{Responsável} & \cab{Prazo} &
\cab{Produto esperado} & \cab{Unidade} & \cab{Evidência} & \cab{Periodic.} &
\cab{Validação} \\ \hline
1. Perfuração de poço &
1.1 Estudo hidrogeológico com locação, projeto construtivo do poço e anuência
prévia; 1.2 Obra de perfuração; 1.3 Teste de vazão; 1.4 Análise
físico-química; 1.5 Outorga; 1.6 Gerenciamento e fiscalização de obras &
IAT / Prefeituras & Ano 0 a 1 &
83 locações e projetos, anuências, testes de vazão, análises físico-químicas e
outorgas & Poços perfurados e outorgados &
Atesto de recebimento e laudos de vistoria & Lote & IAT \\ \hline
2. Instalação do sistema de abastecimento &
2.1 Aquisição de projeto básico e executivo (casa química, reservatório,
eletromecânica e distribuição); 2.2 Aprovação dos projetos; 2.3 Obras de
instalação; 2.4 Projeto \textit{as built}; 2.5 Pré-operação, testes e
capacitação dos operadores locais; 2.6 Gerenciamento e fiscalização das obras &
IAT / Prefeituras & Ano 0 a 1 &
83 projetos aprovados, sistemas de armazenamento e distribuição; XXX
hidrômetros instalados & Hidrômetros instalados &
Atesto de recebimento e laudos de vistoria & Lote & IAT \\ \hline
3. Capacitação dos beneficiários &
3.1 Aquisição de treinamentos sobre operação e manutenção dos sistemas &
IAT / Prefeituras & Ano 4 a 6 & Treinamentos na comunidade ou município &
Pessoas capacitadas & Lista de presença e certificados & Contínua & IAT \\ \hline
4. Supervisão técnica &
4.1 Acompanhamento das instalações e suporte técnico pós-implantação &
IAT / Fiscalização de obras & Ano 2 a 6 & Supervisão e gerenciamento de obras &
Relatório de supervisão & Laudos de vistoria e testes de estanqueidade &
Contínua & IAT \\ \hline
5. Monitoramento e avaliação &
5.1 Avaliação de desempenho dos sistemas, indicadores de impacto e ajustes &
IAT & Ano 2 a 6 & Relatórios de conformidade técnica e ambiental &
Relatório & Dados no SIGARH / GeoPR & Semestral & UGP / MGAS \\ \hline
6. Relatórios e encerramento &
6.1 Consolidação de resultados, relatórios finais e prestação de contas &
IAT & Ano 5 e 6 & Relatório final de gerenciamento e obras concluídas &
Documento final & Termo de Recebimento Definitivo (TRD) & Única &
UGP / Banco Mundial \\ \hline
\end{tabularx}
\fonte{Adaptado de \citeonline[Quadro 29]{mop2026}.}
\end{quadro}
\end{landscape}

\pendencia{Prazos incompatíveis: no \autoref{qua:etapas-3223} (Quadro 29 do
MOP), as etapas 1 e 2 (perfuração e instalação) ocorrem nos anos 0 a 1,
enquanto a ficha-síntese (Quadro 30 do MOP) prevê a atividade nos anos 2 a 5.
Falta também a quantidade de hidrômetros (``XXX'').}

\subsubsection{Interfaces}

\begin{itemize}
  \item \atividade{3.2.2.1}: a entidade gestora que receberá o sistema deve
        estar constituída antes da entrega.
  \item \atividade{3.2.2.2}: define as comunidades e os domicílios atendidos.
  \item \atividade{3.2.2.4}: o esgotamento sanitário será implantado em 70\%
        dos domicílios atendidos por esta atividade.
  \item Salvaguardas: verificação fundiária (MRF) e licenciamento ambiental
        (\autoref{sec:salvaguardas}).
\end{itemize}
